%!TEX root = ../template.tex
%%%%%%%%%%%%%%%%%%%%%%%%%%%%%%%%%%%%%%%%%%%%%%%%%%%%%%%%%%%%%%%%%%%
%% chapter1.tex
%% NOVA thesis document file
%%
%% Chapter with introduction
%%%%%%%%%%%%%%%%%%%%%%%%%%%%%%%%%%%%%%%%%%%%%%%%%%%%%%%%%%%%%%%%%%%

\typeout{NT FILE chapter1.tex}%

\chapter{Introduction}\label{cha:introduction}

The internet has become an integral part of the daily life for most of the world's population. However, in some regions, individuals are denied the ability to freely share their stories or access vital information due to oppressive conditions and restrictions on freedom of speech. This oppression and restrictions are often enforced by totalitarian regimes that have vast resources and power to control the flow of information~\cite*{aryan2013internet, zittrain2017shifting, alimardani2018internet}.

As the demand for anonymity networks continues to expand, so does the necessity for robust privacy guarantees to safeguard secure and confidential communications. At the same time, authoritarian regimes intensify their efforts to restrict access exclusively to state-approved content, employing advanced measures to enforce their censorship policies and maintain control over information flow.

In order to enforce and control the network traffic, these regimes resort to the use of traffic analysis techniques to passively monitor the traffic and infer information about the users. Traffic analysis attacks are a significant threat to anonymity networks, as they can be used to infer users' activities and deny their privacy, as well as deanonymize them~\cite*{chakravarty2014trafficanalysis, winter2012great, robjansen2019dosontor}.
Some have demonstrated that advances in machine learning lead to more sophisticated and effective traffic correlation attacks, such as DeepCorr~\cite{DeepCorr} and DeepCoFFEA~\cite{DeepCoFFEA}, which leverage deep learning techniques to correlate traffic patterns and infer users' activities with high accuracy, in anonymity networks like Tor~\cite{johnson2013users}.
Statistical Disclosure attacks also pose a threat to anonymity networks, as they identify sender-receiver pairs in anonymity networks~\cite*{StatDisclosure,PracticalStatDisclosure}.
Website Fingerprinting attacks represent a major threat to anonymity networks, as they allow attackers to learn information about user's online activities, even on encrypted networks like Tor, by eavesdropping and collecting side-channel information from the network traffic, between the user and the entry node. With this information, the attacker can use machine learning techniques to infer the user's online activities, such as the websites visited, by analyzing the traffic patterns, with high accuracy~\cite{DeepFingerprinting, TikTok, OnlineWebFingerprinting}.

As these attacks become more sophisticated and effective, journalists, whistleblowers and any other person trying to express their freedom of speech and access information in these regions, are at risk of being identified and persecuted, leading to the decrease of this networks utility, mistrust and fear. Therefore, it is critical that anonymity networks become more resilient to these attacks and provide stronger privacy guarantees to their users.


\section{Motivation}\label{sec:motivation}

Tor~\cite{dingledine2004tor}, one of the most popular anonymity networks, is an anonymous communication service based on the Onion Routing Protocol. Tor is designed to protect users' privacy and security by routing their traffic through a network of volunteer-operated servers, encrypting it at each step, and ensuring that no single server knows both the source and destination of the traffic.

%% Intro to Solutions to censorship evasion
Anonymity networks remain vulnerable to attacks such as traffic analysis, which threatens the privacy and security of the users these networks are designed to protect~\cite*{chakravarty2014trafficanalysis, winter2012great, robjansen2019dosontor, StatDisclosure, PracticalStatDisclosure, DeepFingerprinting, TikTok, OnlineWebFingerprinting}. To address this challenge of circumventing censorship, researchers have proposed a range of solutions. Among these solutions, we discuss traffic encapsulation, fully encrypted protocols, and traffic shaping techniques.

%%% Traffic Encapsulation

Traffic encapsulation is a method that conceals data by embedding packets within other data formats or protocols, thereby obfuscating their true content. This technique leverages protocols, such as those for web browsing and video streaming, to disguise network traffic, making it indistinguishable~\cite*{TorKameleon, StegoTorus, Protozoa, Stegozoa, Freewave}.

%%% Fully Encrypted Protocols



%%% Traffic Shapping
Traffic Shaping is another technique that aims to regulate data transmission to optimize network performance and to prevent timing and correlation attacks, contributing for privacy-prevention.  
Similarly, traffic shuffling is a technique that focus on reordering packets, with statistical or cryptographic techniques, to make it difficult for an attacker to correlate the traffic~\cite*{DAENet, Karaoke, Riposte, Loopix, Dissent}.

Researchers have proposed various solutions.~\citeauthor{kAnonymityEffectiveness} demonstrated that K-Anonymity is effective against traffic analysis attacks, such as Statistical Disclosure and long-term intersection attacks~\cite{kAnonymityEffectiveness}.~\citeauthor{nunes2023kfunnels} proposed BriK~\cite{nunes2023kfunnels}, a Tor pluggable transport, that ensures uniform traffic patterns to make users indistinguishable. Some work as Walkie-Talkie~\cite{WalkieTalkie}, Supersequence~\cite{Supersequence} and Glove~\cite{Glove} also use traffic shaping for defending against website fingerprinting attacks.~\citeauthor{TrafficMorphing} proposed Traffic Morphing~\cite{TrafficMorphing}, a traffic shaping technique that uses a Markov Chain to generate traffic patterns that are indistinguishable from the original traffic.


\section{Problem}\label{sec:problem}
Although anonymity networks are designed to provide robust privacy and security, advancements in traffic analysis and correlation techniques have increasingly exposed vulnerabilities, reducing their reliability and trustworthiness for users—particularly in regions where pervasive surveillance by authoritarian regimes is a critical concern. These regimes actively invest in sophisticated surveillance methods and traffic manipulation strategies to maintain control over digital communication.

Tor, despite being one of the most widely used anonymity networks, is not immune to these challenges. It remains vulnerable to several attack vectors, including traffic correlation, congestion-based attacks, timing-based correlations, and fingerprinting techniques.~\cite*{chakravarty2014trafficanalysis, johnson2013users,winter2012great, robjansen2019dosontor} Such vulnerabilities highlight the ongoing arms race between those striving to safeguard online anonymity and those seeking to undermine it through technical and infrastructural advancements.

The absence of standardized privacy metrics for evaluating privacy guarantees in a formal and rigorous manner raises concerns about the effectiveness of anonymizing networks like Tor. Research has demonstrated vulnerabilities in the Tor network to various attacks, including traffic analysis. For instance, \citeauthor{chakravarty2014trafficanalysis} showed that statistical correlation techniques can be used to perform successful traffic analysis attacks on the Tor network. Additionally, \citeauthor{yi2009fingerprinting} demonstrated that user privacy can be compromised through fingerprinting attacks

\section{Goals}\label{sec:goals}

The goal is to add differential privacy (DP) mechanisms into the Tor open-source project to preserve its existing security guarantees while enhancing its privacy protections. By incorporating DP, the aim is to provide quantifiable and adjustable privacy guarantees to users, addressing vulnerabilities to correlation attacks and improving resistance to traffic observability. Furthermore, the project seeks to analyze the trade-offs between privacy and utility within the Tor context and to evaluate the performance of the proposed solution compared to the original Tor network.

As stated in Section, it is important to ensure that anonymity networks are secure against threats like artificial intelligence and machine learning, and to address the lack of standardized privacy metrics to evaluate the performance of anonymity networks.
We aim to create and implement a new algorithm that can be used in anonymity networks that leverages differential privacy to ensure privacy and security with higher performance than existing solutions. We also aim to evaluate the performance of the algorithm and compare it with other anonymity networks to prove that this solution is a viable alternative to usual anonymity networks solutions.
To test the algorithm and its performance, we intend to realize correlation attacks on the existing networks and on our prototype and compare the results. To test performance, we will measure the latency and throughput of the algorithm and compare it with other anonymity networks.


\subsection{Contributions}\label{sub:contributions}
This thesis tries to address the problem of ensuring privacy and security in anonymity networks while maintaining high performance. In particular, the contributions of this thesis are:
\begin{enumerate}
    \item Design a new Differential Privacy Based Algorithm for Privacy Preserved Communication in Anonymity Networks.
    \item Implement the algorithm on top of the Tor network and provide a proof of concept.
    \item Evaluate the performance of the algorithm and compare it with other anonymity networks; This contribution aims to prove if this solution is a viable alternative to usual anonymity networks solutions and if, as mention before, can allow users to switch from traditional to anonymity networks.
\end{enumerate}

\section{Report Organization}\label{sec:report_organization}
The remainder of this report is organized as follows: Chapter introduces anonymity networks and differential privacy, as well as some other related topics. The proposal of work is presented on Chapter, which provides some system model and guidelines for the development and discusses the proof method to verify the work. Finally, workplan is explained and addressed in Chapter